\chapter{Worksheets for six release types}\label{ch:worksheets}

A common view among central bank statisticians is that their release types are too
different for one method. Aggregate tables from a credit register are protected by
contributor and dominance rules, survey microdata by recoding and the journalist
measure, ad hoc requests by a reviewer's judgement, and each tradition has its own
literature \citep{hundepool2012}. On this view a single framework can only be a
lowest common denominator. The view has a fair basis in the tools, which do differ.
It does not hold for the decision. Every one of the six release types in this chapter
fills the same four slots: whom the adversary targets, what it holds, which harm
follows and which measure answers the question. What changes from type to type is
the content of the slots, not their number.

The chapter asks one practical question: what should an institution write in the
worksheets of Chapter~\ref{ch:protocol} for the release types it actually publishes?
It gives starting declarations for six types, then the blank worksheets and a
checklist. The declarations are illustrations to be adapted to each institution's
products, legal basis and users. They do not describe the practice of any
institution.

\section{Starting declarations}

Table~\ref{c17:tab-map} maps six release types common in central banks to the objects
of Steps~2 to~5 of the protocol. The first column names the product. The others say
whom the adversary targets, what it plausibly holds, which harm is at stake, which
measure in this book applies and where a pooled scenario arises. The last column is
the one most often left empty in practice, and it is the one
Proposition~\ref{prop:nopool} says cannot be left empty when records or cells persist
across vintages.

\begin{table}[ht]
\centering
\caption{Starting declarations for six release types common in central banks. The
entries are illustrations to be adapted to each institution's products, legal basis
and users; they do not describe the practice of any institution.}
\label{c17:tab-map}
% source: paper/apply_section.tex tab:apply-map
\footnotesize
\setlength{\tabcolsep}{3pt}
\begin{tabular}{@{}>{\raggedright\arraybackslash}p{2.2cm}>{\raggedright\arraybackslash}p{2.2cm}>{\raggedright\arraybackslash}p{2.5cm}>{\raggedright\arraybackslash}p{2.2cm}>{\raggedright\arraybackslash}p{2.5cm}>{\raggedright\arraybackslash}p{2.4cm}@{}}
\toprule
Release type & Targets & Adversary holds & Harm & Measure & Pooled scenario \\
\midrule
Aggregates from a credit register & Borrowers in small cells; lenders in
concentrated cells & A firm's sector, region and size from public filings; a
lender's own book & Identification of a borrower; a lender's position revealed &
Prosecutor risk by cell; residual disclosure from totals & Monthly vintages;
differencing of consecutive months \\
\addlinespace
Loan-level research extract & Every borrower in the extract & Key fields and a noisy
reading of exposure from other sources & Identification; disclosure of arrears &
The linkage attack of Chapter~\ref{ch:credit}; an attribute rule & Successive
extracts for one project; extracts for different projects \\
\addlinespace
Supervisory statistics by bank group & Individual banks & Each bank's own figures;
published annual reports & A bank's confidential figure recovered & Contribution
from totals; dominance and $p\%$ rules & Quarterly series; tables that share a
margin \\
\addlinespace
Payment and transaction statistics & Merchants or payers in thin cells & Merchant
category and location; own transactions & Identification; commercial information
revealed & Prosecutor risk by cell; minimum contributor count & Daily or weekly
series of the same cells \\
\addlinespace
Household or firm survey microdata & Respondents & Demographic or firm
characteristics from registers or social media & Identification; disclosure of
wealth or debt & Journalist and marketer measures; the frequency model & Survey
waves; a public file and a scientific-use file \\
\addlinespace
Answers to ad hoc requests & Units in the requested breakdown & Everything already
published from the same register & Identification by differencing & Differencing
posterior & Each answer pooled with every earlier publication \\
\bottomrule
\end{tabular}
\end{table}

The six rows are discussed in turn below. Each discussion says what the declaration
must settle before any attack is run, what evidence the earlier chapters offer, and
where the assessment usually goes wrong.

\subsection*{Aggregates from a credit register}

Tables of credit by sector, region and size class are the most frequent product of a
credit register, and they face two adversaries at once. An outsider who knows a
firm's sector, region and size from its public filings looks for the cell where the
firm sits alone or nearly so. A participating lender knows its own contribution to
every cell and can subtract it from a published total to read what remains. The
first is measured by the prosecutor risk of each cell (Chapter~\ref{ch:measures}). The
second is a matter of arithmetic, not probability: on the synthetic register a
participating lender can recover 5.9 per cent of 814 published cells by subtraction,
% source: paper/bis_ifc_full.tex tab:record (5.9 per cent of 814 cells)
which is what the three-contributor rule is written to prevent. A dominance rule
usually sits beside the contributor rule, and the record should say whether it was
assessed. The pooled scenario is the monthly series: a cell that changes by one unit
between two months may reveal a single loan's movement, and the record should carry
the differencing posterior for consecutive months (Chapter~\ref{ch:credit}).

\subsection*{Loan-level research extracts}

A research extract puts every borrower in the file, so the target set is the whole
extract and the adversary holds the key fields plus whatever it can read of the
exposure from other sources. The declaration that decides the result is the
precision of that reading. On the synthetic register, halving its standard deviation
from 0.6 to 0.3 on the log scale moves the programme floor from a median of 150 to a
median of 300 over twenty register draws.
% source: paper/generated/table_replication.tex
The worksheet should therefore record the reading precision and the argument for it.
The pooled scenario arises when one researcher receives several extracts, for one
project over time or for different projects, and the declaration must say whether
the researcher can know in which vintages a target's loan was live. Arrears status
needs an attribute rule as well, since a key class in which every large borrower is
in arrears discloses that status without naming anyone.

\subsection*{Supervisory statistics by bank group}

When a table reports figures by bank group, the targets are individual banks and the
adversary is often another bank, which knows its own figures and can read the
others' published annual reports. The harm is the recovery of a confidential figure,
and the relevant tools are the contributor and dominance rules of the sensitivity
tradition \citep{hundepool2012}. These are rules, not probabilities, and the record
should report them as such, in the way the seventh row of the decision record in
Chapter~\ref{ch:protocol} does. The pooled scenario is the quarterly series and any
pair of tables that share a margin, since a margin shared by two tables can be
subtracted from both.

\subsection*{Payment and transaction statistics}

Payment statistics are published at high frequency, and their thin cells, a merchant
category in a small town, can contain one or two merchants. The adversary knows the
category and the location and, if it is a participant, its own transactions. The
measures are the prosecutor risk by cell and a minimum contributor count. The pooled
scenario is the defining feature of this release type: a daily or weekly series of
the same cells gives an adversary many readings of each, and the horizon of
Step~1 should be set in days or weeks, not in years.

\subsection*{Household or firm survey microdata}

Survey microdata are the case for which most of the classical measures were built.
The targets are respondents, and the adversary holds demographic or firm
characteristics from registers or social media. Because the file is a sample, the
journalist measure, which asks whether a sample-unique record is unique in the
population, needs a model for the population cell sizes, and the transport test of
Chapter~\ref{ch:falsification} applies to that model. The marketer measure, the
expected share of records an adversary would link correctly, answers a different
question and protects a different party, and the worksheet should say which of the
two the tolerance is for. The pooled scenario is a sequence of waves, or a public
file and a scientific-use file cut from the same survey.

\subsection*{Answers to ad hoc data requests}

An ad hoc request is the release type most often assessed alone and least suited to
it. The adversary is anyone who holds everything already published from the same
register, so every answer is pooled with every earlier publication. The harm is
identification by differencing, and the measure is the differencing posterior. The
practical consequence is that the inventory of Step~1 has to include past answers,
which many institutions do not keep in one place.

When Naledi filled the first worksheet for the credit register, the row for ad hoc
answers took the longest. Three departments had answered requests from the register
in the previous year, and two kept no list. She recorded what could be found and
noted the gap in the record, since an inventory that is known to be incomplete is
more useful to the owner than one that looks complete.

\section{Blank worksheets}

The three worksheets below follow the steps of Chapter~\ref{ch:protocol}. They are
meant to be copied into an institution's own templates. The decision record of
Table~\ref{c16:tab-record} shows the third worksheet filled for the synthetic credit
register, with the declarations common to its rows given in the text.

\subsection*{Worksheet 1. Programme inventory (Step 1)}

One row per product drawn from the source register over the planning horizon. The
fifth column is the one that decides whether pooled scenarios are needed.

\begin{center}
\footnotesize
\setlength{\tabcolsep}{3pt}
\begin{tabular}{@{}>{\raggedright\arraybackslash}p{2.6cm}>{\raggedright\arraybackslash}p{1.7cm}>{\raggedright\arraybackslash}p{2.3cm}>{\raggedright\arraybackslash}p{2.6cm}>{\raggedright\arraybackslash}p{3.0cm}>{\raggedright\arraybackslash}p{2.1cm}@{}}
\toprule
Product drawn from the register & Frequency & Population covered & Variables
published & Link across vintages (pseudonym, row order, stable attributes) & Access
route \\
\midrule
 & & & & & \\[12pt]
 & & & & & \\[12pt]
 & & & & & \\[12pt]
 & & & & & \\[12pt]
\bottomrule
\end{tabular}
\end{center}
\noindent Source register: \hrulefill\quad Horizon $m$ (vintages): \hrulefill\quad
Data owner: \hrulefill

\subsection*{Worksheet 2. Scenario declaration (Step 2), one per scenario}

One sheet per declared adversary, single-release or pooled. The stress-test row
names the stronger adversary that the falsification test of
Chapter~\ref{ch:falsification} will run.

\begin{center}
\small
\begin{tabular}{@{}>{\raggedright\arraybackslash}p{5.4cm}>{\raggedright\arraybackslash}p{8.8cm}@{}}
\toprule
Item & Entry \\
\midrule
Scenario name and products it holds & \\[8pt]
Adversary (who, and why it would attack) & \\[8pt]
Auxiliary information held, with its error & \\[8pt]
Target set (prosecutor, journalist or marketer) & \\[8pt]
Success event (what counts as a correct attack) & \\[8pt]
Harm class & \\[8pt]
Single release or pooled over the horizon & \\[8pt]
For a pooled scenario: does the adversary know in which vintages the target is
present, and why & \\[8pt]
Stronger adversary for the stress test & \\[8pt]
Declared by, and date & \\[8pt]
\bottomrule
\end{tabular}
\end{center}

\subsection*{Worksheet 3. Decision record (Steps 3 to 7), one row per scenario}

One row per scenario and harm, including those for which no tolerance has been set.
The decision column takes one of four values: pass, fail, invalid (a falsification
test failed) or not set (no tolerance has been agreed).

\begin{center}
\footnotesize
\setlength{\tabcolsep}{2.5pt}
\begin{tabular}{@{}>{\raggedright\arraybackslash}p{1.5cm}>{\raggedright\arraybackslash}p{1.8cm}>{\raggedright\arraybackslash}p{1.3cm}>{\raggedright\arraybackslash}p{0.6cm}>{\raggedright\arraybackslash}p{0.6cm}>{\raggedright\arraybackslash}p{1.1cm}>{\raggedright\arraybackslash}p{1.5cm}>{\raggedright\arraybackslash}p{1.6cm}>{\raggedright\arraybackslash}p{1.6cm}>{\raggedright\arraybackslash}p{1.4cm}@{}}
\toprule
Scenario & Legal basis and protected parties & Access route & $\tau$ & $\alpha$ &
Attacks $B$ & Upper bound & Model tests passed & Decision & Remedy if
failed \\
\midrule
 & & & & & & & & & \\[12pt]
 & & & & & & & & & \\[12pt]
 & & & & & & & & & \\[12pt]
\bottomrule
\end{tabular}
\end{center}
\noindent Release owner: \hrulefill\quad Signed on: \hrulefill\quad Expiry:
\hrulefill

\noindent Re-assess when the horizon is reached, a new product is drawn from the
register, a new public source could serve as auxiliary information, the mechanism
changes, or the record expires.

\section{Checklist before release}

The checklist is for the reviewer and, afterwards, for internal audit. Each item can
be answered yes or no from the documents.

\begin{enumerate}
\item Every product drawn from the register over the horizon is listed.
\item Every scenario names its adversary, auxiliary information, target set and harm.
\item A tolerance, error level, owner and expiry are recorded for every harm, or the
record says that none has been set.
\item The legal basis, the protected parties and the access route are recorded for
every product.
\item The number of attacks, the error level of each scenario and the order in which
any grid of levels is tested were fixed before any result was seen.
\item Pooled scenarios were run with the same auxiliary information as the
single-release ones; any pooled scenario that adds knowledge is reported separately;
all are reported.
\item The decision uses upper bounds, not point estimates.
\item All six falsification tests were run against bands fixed in advance, and their
results are filed.
\item The reviewer is not the analyst, and the release owner signed the record.
\end{enumerate}

The answer to the chapter's question is the worksheets themselves, filled from the
starting declarations of Table~\ref{c17:tab-map} and changed wherever the
institution's own products, users and law differ from them. The six release types
differ in their targets, adversaries, harms and measures, and they share the four
slots and the pooled column. An institution that fills the slots for its own
products, and leaves none empty without saying why, has done Steps~1 to~3 of the
protocol.

\begin{inpractice}[In practice: using the worksheets]
\setlength{\emergencystretch}{3em}%
Copy the three worksheets into the institution's own templates and keep the column
for links across vintages and the row on the adversary's knowledge of presence, since
those decide whether a pooled scenario is needed. Start each declaration from the
matching row of Table~\ref{c17:tab-map} and record every change from it with a
reason. Keep ad hoc answers in the programme inventory. Record rows for harms that
have no tolerance yet, with the decision ``not set''.

\smallskip\noindent\textit{In the software.} The pipeline fills the decision record
for the declared scenarios of a run (\texttt{privacy\_\allowbreak{}risk.\allowbreak{}pipeline.\allowbreak{}gate}, written to
\texttt{decision.\allowbreak{}json}); the credit register's harms are measured by
\texttt{privacy\_\allowbreak{}risk.\allowbreak{}credit} and attacked by \texttt{privacy\_\allowbreak{}risk.\allowbreak{}evaluation}. The
inventory and scenario sheets have no counterpart in code and are kept as documents.
\end{inpractice}
